\subsubsection{The two exceptional five-dimensional module pairs}

The following calculation is invariant under algebra isomorphism; it does
not assume that an isomorphism preserves the distinguished monoid basis.
All modules in this subsection are finite-dimensional over $\Q$.

\begin{lemma}\label{B:five-exceptional-targets}
Let $A$ be one of the two exceptional algebras in
Proposition~\ref{B:five-nonlocal-exceptions}.  Suppose that a right module $W$
and a left module $V$ satisfy
\[
 \dim V=2,\quad \dim W=4,\quad
 \dim\operatorname{Tor}_1^A(W,V)=0,\quad
 \dim\operatorname{Tor}_2^A(W,V)=1.
\]
For the zero-composite line algebra, the additional condition
$\dim(W\otimes_A V)\leq2$ forces
\[
 V=S_1\oplus S_3,\qquad
 W=T_1\oplus T_3\oplus e_2A,
 \qquad \dim(W\otimes_A V)=2.
\]
Here $S_i$ and $T_i$ denote the left and right simple modules at vertex
$i$.  In particular this algebra cannot give the additive homology of
$\CP^2$ through the band--Tor formula.

For the Nakayama algebra, the conditions force
\[
 V=Ae_1/\Q c,\qquad W=T_2\oplus e_1A,
 \qquad \dim(W\otimes_A V)=1.
\]
The unique, up to scalar, balanced form $W\times V\to\Q$ has matrix
rank one.  Thus this algebra cannot give the additive homology of
$S^2\times S^2$, and its degree-two square vanishes in every rectangular
band realization by Lemma~\ref{B:rank-one-square}.
\end{lemma}

\begin{proof}
For the line algebra, put $a\in e_1Ae_2$, $b\in e_2Ae_3$ and $ab=0$.
A left module is a diagram
$V_3\xrightarrow{b}V_2\xrightarrow{a}V_1$ with zero composite.
Choose a basis of $\operatorname{im}b$, extend it to a basis of
$\ker a$, and then extend to $V_2$; compatible bases in $V_1,V_3$
split the representation into the three simples and the two one-arrow
modules.  The same argument applies on the right.  In particular the
right indecomposables are $T_1,T_2,T_3,e_1A,e_2A$, and their only
nontrivial projective resolutions are
\[
 0\longrightarrow e_3A\xrightarrow{b}e_2A
   \xrightarrow{a}e_1A\longrightarrow T_1\longrightarrow0,
 \qquad
 0\longrightarrow e_3A\xrightarrow{b}e_2A
   \longrightarrow T_2\longrightarrow0.
\]
Consequently $\operatorname{Tor}_2(T_1,V)=\ker(b:V_3\to V_2)$,
$\operatorname{Tor}_1(T_1,V)=\ker a/\operatorname{im}b$, and
$\operatorname{Tor}_1(T_2,V)=\ker b$.

The degree-two Tor condition says that exactly one copy of $T_1$
occurs in $W$ and exactly one copy of $S_3$ occurs in $V$.  Since
$\dim V=2$, the second summand of $V$ must be $S_1$ or $S_2$;
the latter makes $\operatorname{Tor}_1(T_1,V)$ nonzero.  Thus
$V=S_1\oplus S_3$.  The degree-one condition excludes $T_2$ from
$W$, so
\[
 W=T_1\oplus T_3^{\oplus s}\oplus(e_1A)^{\oplus p}
                 \oplus(e_2A)^{\oplus q},\qquad
 s+2p+2q=3.
\]
Its tensor product with $V$ has dimension $1+s+p$.
If this is at most two, parity forces $(s,p,q)=(1,0,1)$.

For the Nakayama algebra, take
$a\in e_1Ae_2$, $b\in e_2Ae_1$, $ba=0$, $ab=c$.
Its left indecomposables are
\[
 S_1, S_2, L=Ae_1/\Q c, Ae_1, Ae_2,
\]
of dimensions $1,1,2,3,2$; the right indecomposables are
$T_1,T_2,R=e_1A/\Q c,e_1A,e_2A$ of the same dimensions.
For completeness, this decomposition is elementary: in a representation
with maps $b:V_1\to V_2$ and $a:V_2\to V_1$, choose a basis of
$\operatorname{im}(ab)$ and lift it successively through $a$ and $b$.
These length-three chains split off copies of $Ae_1$.
On the remaining summand both composites vanish; splitting the two
images inside the opposite kernels gives length-two chains and simples.
The right-handed argument is identical.

The relevant right resolutions are
\[
\begin{split}
 &0\longrightarrow e_2A\xrightarrow{a}e_1A
       \longrightarrow T_1\longrightarrow0,\\
 &0\longrightarrow e_2A\xrightarrow{a}e_1A
       \xrightarrow{b}e_2A\longrightarrow T_2\longrightarrow0,\\
 &0\longrightarrow e_2A\xrightarrow{a}e_1A
       \xrightarrow{c}e_1A\longrightarrow R\longrightarrow0.
\end{split}
\]
Each displayed arrow is specified by the image of the generating
idempotent.  These formulas give the following table; each entry is
$(\dim\operatorname{Tor}_0,\dim\operatorname{Tor}_1,
\dim\operatorname{Tor}_2)$.
\[
\begin{array}{c|ccccc}
 &S_1&S_2&L&Ae_1&Ae_2\\\hline
T_1 &(1,0,0)&(0,1,0)&(1,1,0)&(1,0,0)&(0,0,0)\\
T_2 &(0,1,0)&(1,0,1)&(0,0,1)&(0,0,0)&(1,0,0)\\
R   &(1,1,0)&(0,0,1)&(1,1,1)&(1,0,0)&(1,0,0)\\
e_1A&(1,0,0)&(0,0,0)&(1,0,0)&(2,0,0)&(1,0,0)\\
e_2A&(0,0,0)&(1,0,0)&(1,0,0)&(1,0,0)&(1,0,0)
\end{array}
\]
The five possible dimension-two left modules are
$S_1^{\oplus2},S_1\oplus S_2,S_2^{\oplus2},L,Ae_2$.
The first and last have zero degree-two Tor, the third has even
degree-two Tor dimension, and the second has positive degree-one Tor
whenever its degree-two Tor is nonzero.  Therefore $V=L$.
The table then excludes $T_1,R$ from $W$ and requires exactly one
$T_2$.  The remaining dimension is three, so $W=T_2\oplus e_1A$.

Finally, $T_2\otimes_A L=0$ and $e_1A\otimes_A L=e_1L=\Q$.
In the bases $(e_1,a,c)$ of $e_1A$ and $(\bar e_1,\bar b)$ of $L$,
the balanced pairing is supported only on $(e_1,\bar e_1)$.
It therefore has rank one, as asserted.
\end{proof}

\begin{corollary}[An independent lower bound of twenty-one for
the projective plane]\label{B:CP2-lower21}
Every finite monoid whose classifying space is $\CP^2$ has at
least twenty-one elements.
\end{corollary}
\begin{proof}
Corollary~\ref{B:fourdim-lower20} already excludes orders below
twenty.  At order twenty, apply the same corner and minimal-ideal
reductions.  The two-row/column, three-by-three and three-by-four
obstructions leave a rectangular band with at least fifteen
elements; a nontrivial minimal maximal group would require at
least twenty-four ideal elements.  By
Corollary~\ref{B:five-outside-minimal-band}, at least five elements
lie outside the band.  Thus the only remaining shape, up to
opposite, is $3\times5$, with exactly five outside elements.

Put $A=\Q M/\Q K$.  The band--Tor formula gives modules
$\dim V=2$, $\dim W=4$ with Tor dimensions $1,0,1,0,\ldots$
starting in degree zero.  If $A$ is local, its residue division
algebra has dimension dividing five.  The existence of the
two-dimensional module $V$ rules out residue dimension five,
so Lemma~\ref{B:local-five-targets} applies.  If $A$ is nonlocal,
Proposition~\ref{B:five-nonlocal-exceptions} leaves only its two
exceptional algebras.  Lemma~\ref{B:five-exceptional-targets}
excludes the line algebra by additive homology, and the Nakayama
algebra by its zero degree-two square.  This proves the bound.
\end{proof}
